% ---
% id: meat-core/rem:classes
% kind: rem
% title: "Retention and deletion of sectors"
% summary: >-
%   Retaining every sector computes the value as defined; deleting one on an
%   approximate match computes a restriction whose excess no function of a
%   tolerance bounds, so a tolerance may only refuse a removal. Dropping a
%   constraint-active sector and averaging sectors can both overstate the key.
% prominence: standard
% status: accepted
% uses: [thm:removal, prop:retention, def:classes, prop:removal-compose, prop:no-averaging, prop:value-range, def:feasible]
% mentions: [prop:optol-nobound, rem:target-conditions, rem:removal-sets]
% assumes: []
% refs: []
% provenance: newly-derived
% lean:
% ---

\begin{remark}\label{rem:classes}
\emph{The direction of every removal decision.} Removability of
Definition~\ref{def:classes}\,(d) is a licence to delete, and every outcome
other than all six of (W1)--(W6) --- a nonzero residual, a rounding
operation, an absent witness, a witness that names $s^{*}$, a target that
is not key-inert, a block that is not exactly scalar --- retains the
sector. Retention is unconditionally sound and costs nothing: $\hlow$ is
\emph{defined} by \eqref{eq:hlow-def} as the infimum over the whole of
$\Feas$, and Theorem~\ref{thm:removal} says only that one particular
restriction of that infimum does not change it. A run that retains every
sector computes the definition. A run that deletes one on an
\emph{approximate} match computes instead the infimum over
$\{\omega\in\Feas:\omega_{s^{*}}=0\}$, a restriction, whose value is
$\ge\hlow$ and which would be reported under the name $\hlow$ --- an
overstatement, and one no function of a \tref{declared tolerance} bounds:
the mechanism is a $P_{\hat c}$ or $K_{\hat c}$ driven towards zero inside a
logarithm, which Proposition~\ref{prop:value-range}\,(iii) exhibits driving
the objective up without limit and which
Proposition~\ref{prop:optol-nobound}\,(i) turns, in the neighbouring
setting of a perturbed declared POVM, into two declarations at distance
$\epsilon$ whose values differ by tens of bits per round --- by
$\frac{\alpha}{\alpha-1}N$ with $N$ up to $52$, that ceiling being the
number model's and not the mechanism's. So a tolerance may only ever
\emph{refuse} a removal, never license one.

\emph{Two unsound operations.} Dropping a constraint-active sector
(Proposition~\ref{prop:retention}) and averaging sectors
(Proposition~\ref{prop:no-averaging}) are both unsound, and both failure
modes overstate the key --- the dangerous direction. The first is reachable
in two ways from a checklist that looks complete: by dropping (W2), which
lets a sector witness its own removal (Remark~\ref{rem:target-conditions}),
and by deleting every removable sector at once, which lets two sectors
witness each other's (Remark~\ref{rem:removal-sets}). (W2) refuses the
first; the hypothesis of Proposition~\ref{prop:removal-compose} refuses the
second, and is the reason removability being a property of one sector does
not by itself license the deletion of a set.
\end{remark}
